% =====================================================================
% Semana 1 · Actividades: normas de los resistores
% Fundamentos de Sistemas Electrónicos Analógicos (Ingeniería Aeroespacial)
% Compilar: pdflatex semana_1_actividades_normas_resistores.tex (2 pasadas)
% =====================================================================
\documentclass[aspectratio=169]{beamer}

\usepackage[utf8]{inputenc}
\usepackage[T1]{fontenc}
\usepackage{lmodern}
\usepackage[spanish,es-tabla]{babel}
\usepackage{amsmath,amssymb}
\usepackage{siunitx}
\usepackage{booktabs}
\usepackage{microtype}
\usepackage{circuitikz}

\definecolor{navy}{RGB}{23,54,93}
\definecolor{navysoft}{RGB}{72,101,140}
\definecolor{gold}{RGB}{179,142,62}
\definecolor{ink}{RGB}{44,51,58}
\definecolor{soft}{RGB}{244,241,234}

\usetheme{default}
\setbeamertemplate{navigation symbols}{}
\setbeamertemplate{footline}[frame number]
\setbeamercolor{structure}{fg=navy}
\setbeamercolor{frametitle}{fg=white,bg=navy}
\setbeamercolor{title}{fg=navy}
\setbeamercolor{subtitle}{fg=navysoft}
\setbeamercolor{itemize item}{fg=gold}
\setbeamercolor{itemize subitem}{fg=gold}
\setbeamercolor{block title}{fg=white,bg=navy}
\setbeamercolor{block body}{bg=soft}
\setbeamercolor{example text}{fg=navy}
\setbeamertemplate{itemize item}{\scriptsize\raise1.25pt\hbox{\donotcoloroutermaths$\bullet$}}
\setbeamertemplate{itemize subitem}{\scriptsize\raise1.25pt\hbox{\donotcoloroutermaths$\circ$}}
\setbeamerfont{frametitle}{size=\large,series=\bfseries}

% macro de enlace clicable
\newcommand{\enlace}[2]{\href{#1}{\textcolor{navy}{\underline{#2}}}}

\title[Actividades · Normas de resistores]{Actividades · Normas de los resistores}
\subtitle{Semana 1 · Trabajo autónomo con recursos en línea}
\institute{Licenciatura en Ingeniería Aeroespacial\\ UNAM · ENES Unidad Juriquilla}
\author{M. en C. José de Jesús Santana Ramírez}
\date{Semestre 2027-1}

\begin{document}

% ---------------------------------------------------------------------
\begin{frame}[plain]
  \titlepage
\end{frame}

% ---------------------------------------------------------------------
\begin{frame}{Objetivos y cómo trabajar}
  \begin{itemize}
    \item Leer el \textbf{código de colores} (IEC 60062).
    \item Elegir el \textbf{valor normalizado} (series E12/E24/E96).
    \item Calcular \textbf{tolerancia, peor caso y TCR}.
    \item Seleccionar por \textbf{potencia con derating} (normas aeroespaciales).
  \end{itemize}
  \pause
  \begin{block}{Método por actividad}
    \textbf{¿Qué necesitas saber?} $\rightarrow$ \textbf{Paso a paso} $\rightarrow$
    \textbf{Ejemplo resuelto} $\rightarrow$ \textbf{Tu turno} $\rightarrow$ \textbf{Investiga}
    (clic en el enlace). Toda cantidad con su unidad.
  \end{block}
  \pause
  \begin{block}{Los 4 parámetros que rigen un resistor}
    Valor ($\Omega$) · Tolerancia (\%) · TCR (ppm/°C) · Potencia/derating (W). En
    aeroespacial \emph{deben cumplirse los 4 a la vez}.
  \end{block}
\end{frame}

% ---------------------------------------------------------------------
\begin{frame}{Actividad 1 · Código de colores (IEC 60062)}
  \begin{columns}[T]
    \begin{column}{0.52\textwidth}
      \textbf{Paso a paso}
      \begin{enumerate}
        \item Cuenta las bandas (4, 5 o 6).
        \item Lee los colores de izquierda a derecha.
        \item Dígitos $\times$ multiplicador.
        \item Última banda = tolerancia.
      \end{enumerate}
      \textbf{Ejemplo} \texttt{Marrón-Negro-Rojo-Dorado}: $10\times100 = \SI{1}{\kilo\ohm}\ \pm 5\,\text{\%}$.
    \end{column}
    \begin{column}{0.44\textwidth}
      \footnotesize
      \begin{tabular}{lll}
        \toprule
        Color & Dígito & Mult. \\
        \midrule
        Negro & 0 & $\times1$ \\
        Marrón & 1 & $\times10$ ($\pm1$\,\text{\%}) \\
        Rojo & 2 & $\times100$ \\
        Naranja & 3 & $\times1$k \\
        Amarillo & 4 & $\times10$k \\
        Verde & 5 & $\times100$k \\
        Azul & 6 & $\times1$M \\
        Violeta & 7 & — \\
        Gris & 8 & — \\
        Blanco & 9 & — \\
        Dorado & — & $\times0.1$ ($\pm5$\,\text{\%}) \\
        \bottomrule
      \end{tabular}
    \end{column}
  \end{columns}
\end{frame}

% ---------------------------------------------------------------------
\begin{frame}{Actividad 1 · Tu turno}
  \begin{enumerate}
    \item Marrón–Negro–Rojo–Dorado  $\rightarrow$ \_\_\_\_\_\_ $\Omega\ \pm5\,\text{\%}$
    \item Rojo–Violeta–Marrón–Dorado  $\rightarrow$ \_\_\_\_\_\_ $\Omega\ \pm5\,\text{\%}$
    \item Naranja–Naranja–Naranja–Dorado  $\rightarrow$ \_\_\_\_\_\_ k$\Omega\ \pm5\,\text{\%}$
    \item Marrón–Negro–Negro–Marrón–Marrón  $\rightarrow$ \_\_\_\_\_\_ k$\Omega\ \pm1\,\text{\%}$
  \end{enumerate}
  \pause
  \begin{block}{Investiga (haz clic)}
    \small
    Calculadora de código de colores:
    \enlace{https://www.digikey.com/en/resources/conversion-calculators/conversion-calculator-resistor-color-code}{DigiKey}
    $\;|\;$
    \enlace{https://www.calculator.net/resistor-color-code-calculator.html}{Calculator.net}\\
    Simulador de Ohm para comprobar: \enlace{https://phet.colorado.edu/sims/html/ohms-law/latest/ohms-law_en.html}{PhET Ley de Ohm}.
  \end{block}
\end{frame}

% ---------------------------------------------------------------------
\begin{frame}{Actividad 2 · Series normalizadas (IEC 60063)}
  Se elige el valor de la serie \textbf{más cercano} al cálculo.
  \begin{itemize}
    \item \textbf{E24} (±5\,\text{\%}): \texttt{1.0, 1.1, 1.2, 1.3, 1.5, 1.6, 1.8, 2.0, 2.2, 2.4, 2.7, 3.0, 3.3, 3.6, 3.9, 4.3, 4.7, 5.1, 5.6, 6.2, 6.8, 7.5, 8.2, 9.1}
    \item \textbf{E96} (±1\,\text{\%}): 3 cifras (permite errores menores).
  \end{itemize}
  \vspace{0.3em}
  \textbf{Ejemplo:} $R = \SI{2.37}{\kilo\ohm}$ $\Rightarrow$ E24: \texttt{2.4 k} (error 1.27\,\text{\%}); E96: \texttt{2.37 k} (error 0\,\text{\%}).
  \pause
  \vspace{0.4em}
  \textbf{Tu turno:} $R = \SI{3.10}{\kilo\ohm}$. E24 más cercano $\rightarrow$ \_\_\_\_\_\_ (error \_\_\_\_\_\_\,\text{\%}); ¿existe \texttt{3.09 k} en E96? $\rightarrow$ \_\_\_\_\_\_

  \pause
  \begin{block}{Investiga}
   \small
    Tabla de valores preferidos: \enlace{https://en.wikipedia.org/wiki/E_series_of_preferred_numbers}{Wikipedia E-series}.
  \end{block}
\end{frame}

% ---------------------------------------------------------------------
\begin{frame}{Actividad 3 · Tolerancia y peor caso}
  \[
    R_{min} = R\left(1-\frac{t}{100}\right)
    \qquad
    R_{max} = R\left(1+\frac{t}{100}\right)
  \]
  \textbf{Ejemplo:} $1\,\text{k}\Omega \pm 5\,\text{\%}$ $\Rightarrow$ \SI{950}{\ohm} a \SI{1050}{\ohm}.
  \pause
  \vspace{0.4em}
  \textbf{Tu turno:}
  \begin{enumerate}
    \item $1\,\text{k}\Omega \pm 1\,\text{\%}$ $\rightarrow$ \_\_\_\_\_\_ a \_\_\_\_\_\_ $\Omega$
    \item $4.7\,\text{k}\Omega \pm 1\,\text{\%}$ $\rightarrow$ \_\_\_\_\_\_ a \_\_\_\_\_\_ $\Omega$
    \item Divisor 1:1 con $\pm1\,\text{\%}$ en $V_{in}=5\,\text{V}$: rango de $V_{out}$ y su desvío de 2.5 V.
  \end{enumerate}
  \pause
  \begin{block}{Investiga}
    \small
    Arma y mide el divisor: \enlace{https://phet.colorado.edu/sims/html/circuit-construction-kit-dc/latest/circuit-construction-kit-dc_en.html}{PhET Kit DC}
    $\;|\;$ \enlace{https://www.falstad.com/circuit/}{Falstad}
    $\;|\;$ repasa la simulación \texttt{semana2\_02\_divisor\_tolerancia.cir}.
  \end{block}
\end{frame}

% ---------------------------------------------------------------------
\begin{frame}{Actividad 4 · TCR y deriva térmica}
  \[
    \Delta R = R\cdot TCR\cdot \Delta T
    \qquad
    error\,\text{\%} = 100\cdot\frac{\Delta R}{R}
  \]
  \textbf{Ejemplo:} $R=\SI{10}{\kilo\ohm}$, $TCR=100\,\text{ppm/°C}$, $\Delta T = 50\,^\circ\text{C}$
  $\Rightarrow$ $\Delta R = \SI{50}{\ohm}$ (0.5\,\text{\%}).
  \pause
  \vspace{0.4em}
  \textbf{Tu turno:} con $R = 10\,\text{k}\Omega$, $\Delta T = 50\,^\circ\text{C}$ y
  $TCR = 25\,\text{ppm/°C}$ (película metálica) $\Rightarrow$
  $\Delta R =$ \_\_\_\_\_\_ $\Omega$ (\_\_\_\_\_\_\,\text{\%}). ¿Por qué un divisor de referencia de ADC exige TCR bajo?

  \pause
  \begin{block}{Investiga}
    \small
    \enlace{https://en.wikipedia.org/wiki/Temperature_coefficient}{Wikipedia - Temperature coefficient}
    $\;|\;$ repasa los tipos de resistor y su TCR en las diapositivas de la Clase 2
    (\texttt{docs/presentations/semana\_1\_diapositivas.md}).
  \end{block}
\end{frame}

% ---------------------------------------------------------------------
\begin{frame}{Actividad 5 · Potencia nominal y derating}
  \[
    P = \frac{V^2}{R}
    \qquad
    P_{derated} = 0.5\cdot P_{nominal}
  \]
  \textbf{Ejemplo:} $1\,\text{k}\Omega$, $1/4\,\text{W}$ a \SI{12}{\volt} $\Rightarrow$
  $P=\SI{144}{\milli\watt}$; $P_{derated}=\SI{125}{\milli\watt}$ $\Rightarrow$ \textbf{no cumple} → usar $1/2\,\text{W}$.
  \pause
  \vspace{0.4em}
  \textbf{Tu turno:} $2.2\,\text{k}\Omega$, $1/4\,\text{W}$ a \SI{12}{\volt} $\Rightarrow$
  $P=$ \_\_\_\_\_\_ mW; ¿cumple? $\rightarrow$ \_\_\_\_\_\_

  \pause
  \begin{block}{Investiga}
    \small
    Busca en cada norma la \textbf{tabla de derating de resistores} (no el documento completo):
    \enlace{https://llis.nasa.gov/lesson/676}{NASA LLIS-0676}
    $\;|\;$ \enlace{https://snebulos.mit.edu/projects/reference/NASA-Generic/EEE-INST-002-DeratingPages.pdf}{NASA EEE-INST-002}
    $\;|\;$ \enlace{https://ecss.nl/standard/ecss-q-st-30-11c-derating-eee-components/}{ECSS-Q-ST-30-11C}.
    En la de resistores, la regla es $\approx 50\,\text{\%}$ de la potencia nominal.
  \end{block}
\end{frame}

% ---------------------------------------------------------------------
\begin{frame}{Actividad 6 · Selección integral (ficha)}
  Diseño: divisor que entregue $V_{out}\approx 2.5\,\text{V}$ desde $V_{in}=12\,\text{V}$ (ref. ADC, carga despreciable).
  \[
    \frac{R_2}{R_1+R_2}=0.208 \qquad \text{prueba } R_1=4.7\,\text{k}\Omega,\ R_2\approx \SI{1.24}{\kilo\ohm}\ (\text{E96})
  \]
  \begin{table}
    \small
    \begin{tabular}{lll}
      \toprule
      Parámetro & Qué pedir & Tu decisión \\
      \midrule
      Valor (serie E) & $V_{out}\approx 2.5$ V & R1 = \_\_\_\_, R2 = \_\_\_\_ \\
      Tolerancia & $V_{out}$ en $\pm0.5$\,\text{\%} & \_\_\_\_\,\text{\%} \\
      TCR & deriva $<0.1$\,\text{\%} en $40\,^\circ\text{C}$ & \_\_\_\_ ppm/°C \\
      Potencia & $P\le 0.5\cdot P_{nom}$ & \_\_\_\_ W \\
      \bottomrule
    \end{tabular}
  \end{table}
  \pause
  \begin{block}{Justificación por escrito (vale el 10 \%)}
    Junto a la ficha, redacta \textbf{1–2 líneas por parámetro} explicando \emph{por
    qué} elegiste ese valor, tolerancia, TCR y potencia. La tabla sin razonamiento
    \textbf{no} suma esos puntos.
  \end{block}
  \pause
  \begin{block}{Investiga}
    \small
    Diseño espacial y peor caso: \enlace{https://s3vi.ndc.nasa.gov/ssri-kb/static/resources/electronic-circuit-design-and-analysis-for-space-applications.pdf}{NASA S3VI}
    $\;|\;$ \enlace{https://www.falstad.com/circuit/}{Falstad}.
  \end{block}
\end{frame}

% ---------------------------------------------------------------------
\begin{frame}{Cómo y cuándo entregar}
  \begin{itemize}
    \item \textbf{Formato:} un solo \textbf{PDF} con: (1) respuestas y cálculos
      (puedes escanear tu cuaderno), (2) capturas de las \textbf{simulaciones}
      (PhET/Falstad/LTspice), y (3) la \textbf{ficha + justificación} de la Actividad 6.
    \item \textbf{Archivo:} \texttt{Apellido\_Nombre\_actividad\_resistores.pdf}
    \item \textbf{Plataforma y fecha:} subir a \textbf{aula virtual} antes del
      \textbf{\texttt{fecha límite}} a las \textbf{\texttt{hora}}.
  \end{itemize}
\end{frame}

% ---------------------------------------------------------------------
\begin{frame}{Investigación · 1. Herramientas y estándares}
  \begin{itemize}
    \item \textbf{Código de colores (DigiKey):} comprueba tus cálculos de la
      \textbf{Actividad 1}. ¿Qué información técnica \textbf{adicional} aporta la
      \textbf{6ª banda} en resistores de alta precisión que no existe en los de 4 bandas?
    \item \textbf{Series E (Wikipedia):} localiza la \textbf{ecuación} que define los
      valores de la serie. ¿Por qué la norma \textbf{IEC 60063} usa \textbf{escala
      logarítmica} en lugar de incrementos lineales (subir de 10 $\Omega$ en 10 $\Omega$)?
    \item \textbf{TCR (Wikipedia):} anota los rangos típicos de TCR de tres tecnologías:
      \textbf{carbón}, \textbf{película metálica} y \textbf{alambre enrollado (wirewound)}.
      ¿Cuál es el \textbf{único viable} para el divisor de la Actividad 6?
  \end{itemize}
  \pause
  \begin{block}{Recursos}
    \small
    \enlace{https://www.digikey.com/en/resources/conversion-calculators/conversion-calculator-resistor-color-code}{DigiKey}
    $\;|\;$ \enlace{https://en.wikipedia.org/wiki/E_series_of_preferred_numbers}{Series E}
    $\;|\;$ \enlace{https://en.wikipedia.org/wiki/Temperature_coefficient}{TCR}.
  \end{block}
\end{frame}

% ---------------------------------------------------------------------
\begin{frame}{Investigación · 2. Simulación y verificación}
  \begin{itemize}
    \item \textbf{PhET Ley de Ohm / Kit DC:} arma en el simulador el divisor de la
      \textbf{Actividad 3}, mide $V_{out}$ y anota el resultado. ¿Qué ocurre con la
      \textbf{corriente} si simulas las resistencias en su escenario de \textbf{peor caso}?
    \item \textbf{Falstad:} dibuja el circuito final de la \textbf{Actividad 6}. Al pasar
      el cursor sobre las resistencias, la herramienta muestra la \textbf{potencia}
      disipada en tiempo real. Captura el dato y verifica si cumple tu margen de derating.
  \end{itemize}
  \pause
  \begin{block}{Recursos}
    \small
    \enlace{https://phet.colorado.edu/sims/html/ohms-law/latest/ohms-law_en.html}{PhET Ley de Ohm}
    $\;|\;$ \enlace{https://phet.colorado.edu/sims/html/circuit-construction-kit-dc/latest/circuit-construction-kit-dc_en.html}{PhET Kit DC}
    $\;|\;$ \enlace{https://www.falstad.com/circuit/}{Falstad}.
  \end{block}
\end{frame}

% ---------------------------------------------------------------------
\begin{frame}{Investigación · 3. Normativas aeroespaciales}
  \begin{itemize}
    \item \textbf{NASA LLIS-0676 (Lessons Learned):} describe en un \textbf{párrafo} el
      mecanismo principal de \textbf{falla térmica} de un componente cuando opera en el
      \textbf{vacío} del espacio, frente a su disipación en la atmósfera terrestre.
    \item \textbf{NASA EEE-INST-002 / ECSS-Q-ST-30-11C:} navega al \textbf{índice} y busca
      la \textbf{tabla de derating para resistores}. Extrae las \textbf{dos reglas de oro}:
      \begin{itemize}
        \item ¿Cuál es el \textbf{\% máximo} de potencia nominal (factor de derating) permitido?
        \item ¿Cuál es la \textbf{temperatura máxima} permitida para el componente?
      \end{itemize}
    \item \textbf{NASA S3VI:} enumera \textbf{2 riesgos} principales de usar resistores
      comerciales (COTS) sin clasificar en lugar de componentes con grado espacial
      (ej. radiación, ciclado térmico).
  \end{itemize}
  \pause
  \begin{block}{Recursos}
    \small
    \enlace{https://llis.nasa.gov/lesson/676}{NASA LLIS-0676}
    $\;|\;$ \enlace{https://snebulos.mit.edu/projects/reference/NASA-Generic/EEE-INST-002-DeratingPages.pdf}{NASA EEE-INST-002}
    $\;|\;$ \enlace{https://ecss.nl/standard/ecss-q-st-30-11c-derating-eee-components/}{ECSS-Q-ST-30-11C}
    $\;|\;$ \enlace{https://s3vi.ndc.nasa.gov/ssri-kb/static/resources/electronic-circuit-design-and-analysis-for-space-applications.pdf}{NASA S3VI}.
  \end{block}
\end{frame}

% ---------------------------------------------------------------------
\begin{frame}{Rúbrica y recursos}
  \begin{columns}[T]
    \begin{column}{0.44\textwidth}
      \begin{table}
        \small
        \begin{tabular}{lr}
          \toprule
          Criterio & Peso \\
          \midrule
          Código de colores & 15\,\text{\%} \\
          Serie E & 20\,\text{\%} \\
          Tolerancia/peor caso & 20\,\text{\%} \\
          TCR & 15\,\text{\%} \\
          Potencia/derating & 20\,\text{\%} \\
          Ficha justificada & 10\,\text{\%} \\
          \bottomrule
        \end{tabular}
      \end{table}
    \end{column}
    \begin{column}{0.52\textwidth}
      \small\textbf{Recursos en línea}
      \begin{itemize}
        \item \enlace{https://www.digikey.com/en/resources/conversion-calculators/conversion-calculator-resistor-color-code}{Código de colores (DigiKey)}
        \item \enlace{https://en.wikipedia.org/wiki/E_series_of_preferred_numbers}{Series E (Wikipedia)}
        \item \enlace{https://phet.colorado.edu/sims/html/ohms-law/latest/ohms-law_en.html}{PhET Ley de Ohm}
        \item \enlace{https://phet.colorado.edu/sims/html/circuit-construction-kit-dc/latest/circuit-construction-kit-dc_en.html}{PhET Kit DC}
        \item \enlace{https://www.falstad.com/circuit/}{Falstad}
        \item \enlace{https://llis.nasa.gov/lesson/676}{NASA LLIS-0676}
        \item \enlace{https://snebulos.mit.edu/projects/reference/NASA-Generic/EEE-INST-002-DeratingPages.pdf}{NASA EEE-INST-002}
        \item \enlace{https://ecss.nl/standard/ecss-q-st-30-11c-derating-eee-components/}{ECSS-Q-ST-30-11C}
        \item \enlace{https://s3vi.ndc.nasa.gov/ssri-kb/static/resources/electronic-circuit-design-and-analysis-for-space-applications.pdf}{NASA S3VI}
        \item \enlace{https://en.wikipedia.org/wiki/Temperature_coefficient}{TCR (Wikipedia)}
      \end{itemize}
    \end{column}
  \end{columns}
\end{frame}

\end{document}
